\section{Internally chosen compatible analytic vertices and coordinate Bennett}

Let $R=\mathcal O_{\mathbb C^{d+t},0}$, let $I\subset R$ be an
original analytic ideal, and write $T$ for the actual unshifted Taylor
map. The tangent centre is $\mathbb C^d\times\{0\}$, with actual
normal ideal $N$. Let $V$ be the true degree-lex vertices of $T(I)$,
let $F_a$ be the canonical reduced formal member at $a$, and put
$w_a=\sum_{i=d}^{d+t-1}a_i$. The hypothesis in this chapter is the
literal whole formal centre-power property
\[
 F_a\in\widehat N^{w_a}\qquad(a\in V).
\]
It is imposed on the canonical reduced members of the original ideal.
Deriving it from constancy of the original Hilbert--Samuel function
along a general smooth centre is a separate assertion.

\begin{theorem}[Compatible analytic vertex coefficient data are internally chosen]
\label{hir:actual-reduced-centre-compatible-analytic-vertex-data}
\uses{hir:actual-analytic-centre-compatible-monic-vertex-family,hir:actual-analytic-vertex-coefficient-data,hir:actual-finite-local-coherent-coefficient-ideals,hir:actual-formal-normal-power-contracts-exact-original-analytic-power}
\lean{Hironaka.HironakaActualReducedCentreBennett.exists_actual_reduced_centre_vertex_data}
\leanok
Under the displayed formal centre-power hypothesis there exist genuine
analytic vertex coefficient data $P$ for the same original $I$ such
that every actual generator satisfies $P_a\in N^{w_a}$. In particular
these are internally chosen monic generators of the entire original
ideal, with the actual initial exponent $a$, one genuine common tangent
domain for their finitely many required low coefficients, and their
entire original coefficient Taylor identities. No analytic family or
coefficient presentation is supplied as a hypothesis.
\end{theorem}
\begin{proof}
Apply the analytic centre-vertex construction to the original ideal
and the actual normal coordinate set. It constructs actual germs
$g_a\in I\cap C^{w_a}$ with initial exponent $a$, coefficient one at
$a$, and whole span $I$, where $C$ is the analytic coordinate ideal.
Their entire Taylor series belong to the corresponding whole formal
normal powers, by the genuine coordinate Taylor map. Whole formal
normal-power contraction therefore gives $g_a\in N^{w_a}$.

The actual vertex set is finite. Put $q=\sum_{a\in V}w_a$, so that
$w_a\le q$ for each vertex. Apply the finite actual coefficient-ideal
construction to this same family $(g_a)$ at cutoff $q$. It chooses a
single real open tangent neighbourhood containing zero and actual
coefficient representatives whose entire Taylor series are the literal
normal coefficients of $T(g_a)$. All coefficients of degree less than
$q$ are holomorphic on this neighbourhood, generate genuine coherent
principal germ ideals there, and belong to their original stalks.
These conclusions hold in particular in degree less than $w_a$.
Retain the original monicity, initial exponents, membership and whole
generation of the same $g_a$. They give the stated genuine vertex data
$P$, while their already proved normal-power membership is retained.
The construction also covers an empty vertex set and weights zero.
\end{proof}

\begin{theorem}[All normalized HS orders decrease from canonical reduced centre powers]
\label{hir:actual-reduced-centre-normalized-bennett}
\uses{hir:actual-reduced-centre-compatible-analytic-vertex-data,hir:actual-normal-normalized-hs-le-translated-components,hir:actual-every-translated-normal-vertex-hs-bound}
\lean{Hironaka.HironakaActualReducedCentreBennett.hilbertSamuelValue_actual_reduced_centre_normalized_le}
\leanok
Under the same literal formal centre-power hypothesis, for every normal
pivot $j$, every actual chart parameter $b$, and every order $\ell$,
the original HS value of the actual normalized-member ideal is at most
the original HS value of $I$. No compatible analytic presentation is
an input.
\end{theorem}
\begin{proof}
Choose the genuine data $P$ by the preceding theorem, retaining
$P_a\in I\cap N^{w_a}$. The actual exceptional quotient and zero-axis
restriction theorem therefore bounds the HS of the actual normalized
ideal by the HS of the entire translated normal-symbol ideal of this
same $P$. Here the translation has pivot entry one and its other
entries are the actual normal entries of $b$; tangent variables remain
unchanged. This comparison uses the actual normal-unit coordinate
equivalence and both directions of its whole ideal dictionary. The
translated-vertex theorem for every genuine analytic vertex family
bounds that last HS value by the original HS of the same $I$. Compose
the two inequalities at the arbitrary order $\ell$.
\end{proof}

\begin{theorem}[Original geometric coordinate Bennett from canonical reduced centre powers]
\label{hir:actual-reduced-centre-geometric-bennett}
\uses{hir:actual-reduced-centre-normalized-bennett,hir:actual-normal-geometric-strict-normalized-inclusion,hir:actual-normal-blowup-codimension-one-identity}
\lean{Hironaka.HironakaActualReducedCentreBennett.hilbertSamuelAt_actual_reduced_centre_strict_le}
\leanok
Let $W\subset\mathbb C^{d+t}$ be any subset, and let $I$ be its
entire original vanishing germ ideal at zero. Suppose the canonical
reduced formal members of this $I$ satisfy the displayed whole normal
power condition. Then, for every normal pivot, chart parameter and
order, the original intrinsic HS of the original geometric strict
transform at zero is at most the original intrinsic HS of $W$ there.
\end{theorem}
\begin{proof}
If $t=1$, the actual normal blowup model is the identity on the whole
ambient space and has empty geometric exceptional locus. Its actual
strict transform is the closure of $W$, whose entire original
vanishing ideal is the same as that of $W$. Thus all intrinsic HS
values are equal. Otherwise an actual normal pivot gives $t>0$, hence
$t\ge2$. The genuine exceptional-locus theorem and the original
exceptional product identity give whole normalized-ideal inclusion in
the vanishing ideal of the geometric strict transform. Original HS
antitonicity bounds its intrinsic HS by the normalized HS. Apply the
preceding theorem to $I$ and use the original model-zero whole-stalk
dictionary to identify the remaining HS value with the intrinsic one
of $W$. Normal codimension zero has no pivot.
\end{proof}
